\documentclass[11pt]{article}

\usepackage{amsmath, amssymb}
\usepackage{algorithm}
\usepackage{algpseudocode}
\usepackage{geometry}
\geometry{margin=1in}

\title{Blueboard: mathematical background}
\author{}
\date{}

\begin{document}
\maketitle

\section*{Purpose}
This note derives the mathematics behind the non-trivial algorithms referenced in
\texttt{PLAN.md}: representing pen strokes as vector curves, fitting freehand input to
those curves, flattening curves for GPU rendering, simplifying curves for level-of-detail
(LOD) rendering, and the spatial index used for viewport culling. Each section states the
problem, the method chosen, and why.

\section{Bézier curves}

\subsection{Definition}
A cubic Bézier curve is defined by four control points
$P_0, P_1, P_2, P_3 \in \mathbb{R}^2$ and parameter $t \in [0,1]$:
\begin{equation}
  B(t) = (1-t)^3 P_0 + 3(1-t)^2 t\, P_1 + 3(1-t) t^2 P_2 + t^3 P_3.
\end{equation}
Equivalently, in matrix form with $\mathbf{t} = \begin{pmatrix}t^3 & t^2 & t & 1\end{pmatrix}$:
\begin{equation}
  B(t) = \mathbf{t}
  \begin{pmatrix}
    -1 & 3 & -3 & 1 \\
     3 & -6 & 3 & 0 \\
    -3 & 3 & 0 & 0 \\
     1 & 0 & 0 & 0
  \end{pmatrix}
  \begin{pmatrix} P_0 \\ P_1 \\ P_2 \\ P_3 \end{pmatrix}.
\end{equation}
$P_0, P_3$ are the curve's endpoints; $P_1, P_2$ pull the curve's tangent direction and
sharpness at those endpoints without the curve passing through them.

A stroke in Blueboard is stored as a path: a sequence of such cubic segments sharing
endpoints, $\{B_1, B_2, \dots, B_n\}$ with $B_i(1) = B_{i+1}(0)$. This is exact —
no sampling error — and is what makes the stroke renderable losslessly at any zoom level
(the brief's 1\%--200\% requirement): the stored representation has no notion of pixel
resolution at all.

\subsection{De Casteljau's algorithm}
Evaluating $B(t)$ directly from the matrix form is numerically fine, but De Casteljau's
algorithm is the standard approach because it is numerically stable and gives subdivision
``for free'' (needed in \S\ref{sec:flatten}). It recursively lerps between control points:
\begin{align}
  P_i^{(0)} &= P_i, & i = 0,1,2,3 \\
  P_i^{(k)} &= (1-t)\,P_i^{(k-1)} + t\,P_{i+1}^{(k-1)}, & k = 1,2,3,\ i = 0,\dots,3-k
\end{align}
and $B(t) = P_0^{(3)}$. A side effect: the intermediate points
$P_0^{(0)}, P_0^{(1)}, P_0^{(2)}, P_0^{(3)}$ and
$P_0^{(3)}, P_1^{(2)}, P_2^{(1)}, P_3^{(0)}$ are exactly the control points of the two
cubic Béziers obtained by splitting the original curve at $t$. This is how adaptive
subdivision (\S\ref{sec:flatten}) splits a curve without re-deriving new control points
from scratch.

\section{Fitting freehand input to a Bézier path}

Raw pen input is a sequence of sampled points $Q_0, \dots, Q_m$ (from the pointer device),
not curves. Storing them as a polyline would defeat the whole point of vector strokes
(\S1), so they are fit to a small number of cubic Bézier segments. Blueboard uses the
standard method for this, due to Schneider (used in e.g. FontForge/Inkscape's own stroke
fitting):

\begin{enumerate}
  \item \textbf{Parameterize} the sampled points by estimated arc length, giving each
    $Q_j$ a $t_j \in [0,1]$ (chord-length parameterization).
  \item \textbf{Fit} one cubic Bézier by least squares: fix endpoint tangent directions from
    the first/last few samples (so the curve's direction matches the pen's motion at the
    stroke's ends), and solve for the two free scalars (how far $P_1, P_2$ sit along those
    tangent directions) that minimize
    $\sum_j \lVert B(t_j) - Q_j \rVert^2$.
  \item \textbf{Measure error}: the maximum distance from any $Q_j$ to the fitted curve.
  \item \textbf{Subdivide if needed}: if the error exceeds a tolerance $\varepsilon$ (a few
    screen pixels at the zoom level the stroke was drawn at), split the point sequence at
    the point of worst error and recurse on each half, producing an additional curve
    segment there.
\end{enumerate}
This is the classic bias/variance tradeoff made explicit: more segments always fit better,
so $\varepsilon$ is the knob that keeps "smooth natural pen curve" from degenerating into
"one Bézier segment per input sample" (which would be the vector equivalent of just storing
the polyline).

\section{Flattening: curves to GPU triangles}
\label{sec:flatten}

GPUs don't render Bézier curves directly — \texttt{lyon} tessellates each stroke's outline
into triangles. The first step of that is \emph{flattening}: approximating the curve by a
polyline accurate enough that, on screen, no viewer can tell the difference.

\subsection{Flatness test}
A cubic segment is "flat enough" to replace with the single line segment $P_0 P_3$ when its
control points are close enough to that line. A common test bounds the squared distance of
$P_1, P_2$ from line $P_0P_3$ against a tolerance $\varepsilon^2$ (in \emph{screen-space}
pixels, so the tolerance automatically tightens as the user zooms in):
\begin{equation}
  d(P_1)^2 \le \varepsilon^2 \quad \text{and} \quad d(P_2)^2 \le \varepsilon^2,
\end{equation}
where $d(P)$ is the perpendicular distance from $P$ to line $P_0P_3$.

\subsection{Adaptive subdivision}
\begin{algorithm}[h]
\caption{Flatten(curve, $\varepsilon$)}
\begin{algorithmic}[1]
\If{curve is flat enough (flatness test, tolerance $\varepsilon$)}
  \State emit line segment $P_0 \to P_3$
\Else
  \State $(\text{left}, \text{right}) \gets$ De Casteljau split at $t=0.5$
  \State Flatten(left, $\varepsilon$)
  \State Flatten(right, $\varepsilon$)
\EndIf
\end{algorithmic}
\end{algorithm}
Because the screen-space tolerance $\varepsilon$ is fixed in pixels, a curve zoomed in
further gets subdivided more finely automatically — this is the mechanism, not LOD
simplification (\S\ref{sec:lod}), that keeps curves looking smooth at 200\% zoom without
wasting triangles at 1\% zoom.

\section{Level-of-detail simplification}
\label{sec:lod}

\S\ref{sec:flatten} controls rendering fidelity of a single curve at its current zoom, but
at very low zoom (many strokes, each tiny on screen), tessellating full detail for every
stroke is wasted GPU/CPU work even with a loose $\varepsilon$ — the fixed per-stroke
overhead (De Casteljau recursion, buffer uploads) dominates. For this, Blueboard
pre-simplifies a stroke's \emph{flattened polyline} (not the Bézier control points) using
the Ramer--Douglas--Peucker algorithm when a stroke's on-screen bounding box falls below a
pixel-size threshold:

\begin{algorithm}[h]
\caption{DouglasPeucker(points $[Q_0,\dots,Q_n]$, $\varepsilon$)}
\begin{algorithmic}[1]
\If{$n < 2$}
  \State \textbf{return} $[Q_0, \dots, Q_n]$
\EndIf
\State find $Q_k$ with maximum perpendicular distance $d_{\max}$ from line $Q_0 Q_n$
\If{$d_{\max} \le \varepsilon$}
  \State \textbf{return} $[Q_0, Q_n]$
\Else
  \State \textbf{return} DouglasPeucker($[Q_0,\dots,Q_k]$, $\varepsilon$) $\;\cup\;$
         DouglasPeucker($[Q_k,\dots,Q_n]$, $\varepsilon$)
\EndIf
\end{algorithmic}
\end{algorithm}

The simplified polyline is cached per LOD bucket (a small set of zoom-level ranges) alongside
the full-detail tessellation from \S\ref{sec:flatten}, so switching zoom bands is a cache
lookup, not a recomputation (PLAN.md Phase 6).

\section{Spatial indexing: R-tree viewport culling}

With potentially thousands of elements on a semi-infinite board, testing every element
against the viewport every frame is $O(n)$ per frame — the thing that makes "infinite
canvas" expensive if done naively. Blueboard indexes element bounding boxes in an R-tree.

\subsection{Structure}
An R-tree is a height-balanced tree where every node stores the minimum bounding rectangle
(MBR) of everything beneath it; leaves store individual elements' MBRs. Insertion picks the
child subtree whose MBR would need to grow least to contain the new element (minimizing
wasted area), splitting a node into two when it overflows a fixed fanout $M$.

\subsection{Query}
A viewport query (``give me every element whose MBR intersects rectangle $V$'') descends the
tree, at each node skipping any child whose MBR does not intersect $V$, and only recursing
into children that do:
\begin{algorithm}[h]
\caption{Query(node, $V$)}
\begin{algorithmic}[1]
\If{node is a leaf}
  \State \textbf{return} \{elements in node whose MBR intersects $V$\}
\Else
  \State result $\gets \emptyset$
  \For{each child $c$ of node}
    \If{MBR($c$) intersects $V$}
      \State result $\gets$ result $\cup$ Query($c$, $V$)
    \EndIf
  \EndFor
  \State \textbf{return} result
\EndIf
\end{algorithmic}
\end{algorithm}
Expected query cost is $O(\log n + k)$ for $n$ total elements and $k$ results actually
intersecting the viewport, versus $O(n)$ for a linear scan — this is the asymptotic
difference that makes panning around a huge board stay cheap regardless of how much content
exists far outside the current view. (Implementation: the \texttt{rstar} crate, rather than
a hand-rolled tree — same reasoning as using \texttt{lyon} for tessellation rather than
reimplementing it: the algorithm is standard and worth getting from a maintained library.)

\end{document}
